\documentclass[10pt]{article}
\usepackage[T1]{fontenc}
\usepackage{lmodern}
\usepackage[margin=0.75in]{geometry}
\usepackage{amsmath,amssymb,amsthm,booktabs,array,microtype}
\usepackage[hidelinks]{hyperref}
\setlength{\parindent}{0pt}
\setlength{\parskip}{2.5pt}
\newcommand{\Z}{\mathbb Z}
\newcommand{\Kh}{\operatorname{Kh}}
\newcommand{\eps}{\epsilon}
\newtheorem{theorem}{Theorem}
\title{A crossing change alters two integral Khovanov groups\\\large Disproof of the alternating clause of Problem 00000000620}
\author{lizaixi01}
\date{October 24, 2026}
\begin{document}
\maketitle
\vspace{-1.5em}

\textbf{Scope.} The printed conjecture says that every crossing change of an alternating
knot perturbs exactly one Khovanov group. We use the standard unreduced,
normalized integral groups $\Kh^{i,j}$ of the graded $\Z$-module cube in
Bar-Natan's handout~\cite{BN}. The closed positive two-braid $D=\widehat{\sigma_1^3}$
is a one-component alternating diagram. Change its last crossing to obtain
$D'=\widehat{\sigma_1\sigma_1\sigma_1^{-1}}$. We compute both actual cubes directly.
\begin{theorem}
The groups $\Kh^{2,5}(D)$ and $\Kh^{3,9}(D)$ are nonzero, whereas
$\Kh^{2,5}(D')=\Kh^{3,9}(D')=0$. Thus this single crossing change alters at least
two distinct graded groups.
\end{theorem}

\textbf{The actual diagrams and smoothings.}
Number the crossings $r=0,1,2$ from top to bottom. At each crossing the four
ports $4r+p$, $p=0,1,2,3$, are top-left, top-right, bottom-left, bottom-right.
The braid strands pair top-left with bottom-right and top-right with bottom-left.
Closure arcs join each bottom-left/right to the next top-left/right, cyclically.
At a positive crossing the top-left strand passes over the top-right strand;
only this choice at $r=2$ is reversed in $D'$.

An oriented entrance is $(r,b)\in\Z/3\Z\times\{0,1\}$, and the actual next entrance
is $(r+1,1-b)$. This permutation is one six-cycle. For $D$ the over/under
status changes at every successor, including the closure. Hence the diagram is
a knot and is alternating. Both assertions, the correspondence with the port
pairings, and the unique change at $r=2$ are checked in Lean.

Let $s=(s_0,s_1,s_2)\in\{0,1\}^3$. At a positive crossing, smoothing $0$ is vertical
and smoothing $1$ joins the two top ports and the two bottom ports horizontally.
At the negative crossing these two choices are interchanged. A smoothed circle
is an actual connected component of the graph containing the closure pairings
and these smoothing pairings. Its representative below is its least port.
\begin{center}
\begin{tabular}{c@{\quad}c@{\quad}c@{\qquad}c@{\quad}c@{\quad}c}
\toprule
$s$ & circle representatives for $D$ & for $D'$ & $s$ & for $D$ & for $D'$\\
\midrule
000 & $0,1$ & $0$ & 100 & $0$ & $0,2$\\
001 & $0$ & $0,1$ & 101 & $0,2$ & $0$\\
010 & $0$ & $0,6$ & 110 & $0,2$ & $0,2,6$\\
011 & $0,6$ & $0$ & 111 & $0,2,6$ & $0,2$\\
\bottomrule
\end{tabular}
\end{center}
These are verified partitions: every graph edge preserves its representative,
and each port has a concrete graph path to its representative. The resulting
circle type is equivalent to the native graph connected-component quotient.

\textbf{The full integer cube.}
Use $A=\Z[\eps]/(\eps^2)$ with $v_+=1$, $v_-=\eps$, of quantum degrees $1,-1$.
Each state contributes $\bigotimes_{\text{circles}} A$. The saddles use the actual
multiplication and comultiplication
\[
 m(1\otimes1)=1,\quad m(1\otimes\eps)=m(\eps\otimes1)=\eps,\quad
 m(\eps\otimes\eps)=0,
\]
\[
 \Delta(1)=1\otimes\eps+\eps\otimes1,\qquad
 \Delta(\eps)=\eps\otimes\eps.
\]
Unaffected circles are matched by their identical port components; the tensor
factor reindexing is an actual equivalence. An edge changing $s_r=0$ to $1$
has sign $(-1)^{\sum_{k<r}s_k}$. If a decoration has labels $\ell_c\in\{0,1\}$,
its two normalized degrees are
\[
 i=|s|-n_-,\qquad j=\sum_c(1-2\ell_c)+|s|+n_+-2n_-.
\]
Here $(n_+,n_-)=(3,0)$ for $D$, and $(2,1)$ for $D'$. The complete cube differential
preserves $j$, raises $i$ by one, and squares to zero. All these identities are
verified on every actual enhanced state. The enhanced-state counts by
homological degree are $4,6,12,8$ in degrees $0,1,2,3$ for $D$ and
$2,12,12,4$ in degrees $-1,0,1,2$ for $D'$.

\newpage
\textbf{Two non-boundary classes for $D$.}
At $(i,j)=(2,5)$ order the six basis vectors as
\[
 011:(+,-),(-,+);\qquad101:(+,-),(-,+);\qquad110:(+,-),(-,+),
\]
with circle factors in increasing representative order. The incoming degree-one
states are $001,010,100$, each labelled $+$; in degree three order the three
decorations of $111$ with exactly one $+$ by the position of that label.
The actual incoming and outgoing matrices are
\[
 B=\begin{pmatrix}
 1&-1&0\\1&-1&0\\1&0&-1\\1&0&-1\\0&1&-1\\0&1&-1
 \end{pmatrix},\qquad
 C=\begin{pmatrix}
 1&0&-1&0&1&0\\1&0&0&-1&0&1\\0&1&0&-1&1&0
 \end{pmatrix}.
\]
For $z=(0,-1,1,0,1,0)^t$ and $\lambda=(-1,1,0,0,0,0)$,
\[
 Cz=0,\qquad \lambda B=0,\qquad \lambda z=-1.
\]
Thus $z$ is an integer cycle outside the integer image of $B$; it gives a
nonzero class in $\Kh^{2,5}(D)$.

At $(3,9)$ take the $111$ state with all three circles labelled $+$. Its
outgoing differential is zero. Every incoming split is a $\Delta$ and introduces
at least one $-$, so the coordinate functional of this all-$+$ vector vanishes
on all boundaries and has value $1$ on the vector. Its class in
$\Kh^{3,9}(D)$ is therefore nonzero. Lean verifies both probes on the full differential.

\textbf{The changed diagram has zero groups in both places.}
For $D'$, the top degree $i=2$ has only state $111$, with two circles, represented
by $0,2$. Let $t_{00},t_{01},t_{10},t_{11}$ be its four labelings, with $0=+$
and $1=-$. In degree one let
\[
 p_{000}=110:(+,+,+),\quad p_{001}=110:(+,+,-),\quad
 p_+=011:(+),\quad p_-=011:(-).
\]
The circle order for $110$ is $0,2,6$. Direct evaluation of the genuine saddles gives
\[
 d p_{000}=t_{00},\qquad d(p_+-p_{001})=t_{01},\qquad
 d p_{001}=t_{10},\qquad d p_-=t_{11}.
\]
These are integer preimages and preserve quantum degree: the four target degrees
are $5,3,3,1$. Thus the final differential is surjective in every quantum degree,
giving $\Kh^{2,j}(D')=0$ for every $j$. There are no enhanced states in degree
$3$, so also $\Kh^{3,j}(D')=0$. This proves the theorem.

\textbf{Formal objects and identification boundary.}
The Lean project constructs the exact port pairings, actual oriented traversal,
actual smoothing graph components, native dual-number algebra and tensor-product
bases, full signed and doubly graded integer differential, and the groups
\[
 \Kh^{i,j}=\bigl(C^{i,j}\cap\ker d\bigr)/
       \operatorname{range}\bigl(d:C^{i-1,j}\longrightarrow C^{i,j}\cap\ker d\bigr)
\]
as native submodules and quotient modules. The two nonzero quotient classes and
the changed groups' subsingleton properties exclude additive group isomorphisms.
The finite representative tables are certified against actual graph paths.

The identification with the standard unreduced Khovanov construction uses the
published diagram/cube definition in~\cite{BN}. Ambient isotopy of spatial
embeddings and the general Reidemeister-invariance theorem are outside this
Lean development. Neither trefoil classification nor recognition of $D'$ as an
unknot is used in the computation. The conclusion addresses only the printed
alternating-knot every-crossing clause; it makes no claim about the separate
nonalternating parenthesis.

\textbf{Reproduction.} The pinned project is under \texttt{lean/}. Run
\texttt{lean --version}, then \texttt{lake build +Main},
\texttt{lake env lean -DwarningAsError=true Check.lean}, and
\texttt{lake env leanchecker --fresh Main}. The evidence records the exact source,
all named declarations and generated names, their ordinary axiom audit, and the
source-bound execution. Production timings carry no controlled efficiency claim.

\begin{thebibliography}{1}\small
\bibitem{BN} D.~Bar-Natan, \emph{Khovanov Homology For Knots},
2022--23 Fast Computations in Knot Theory,
\href{https://katlas.math.toronto.edu/drorbn/AcademicPensieve/Classes/23-FastComputations/KH4Knots/KH4Knots.pdf}{author's complete one-page handout}.
Its $\Z$-module definition, $m,\Delta$, shifts, signs and eight trefoil resolutions are used above.
\end{thebibliography}
\end{document}
